\scanpage{37}
On peut encore dire que l'on a un
accouplement \uuline{alterné}
\[
 V\otimes V\xrightarrow{\varphi}\deleted{\unsure{K}}L \tag*{(57 bis)}
\]
défini par la condition que l'on a
\[
\varphi(i_1(\underset{\substack{\rotatebox{-90}{$\in$}\\\Gamma L}}{\lambda})\otimes\underset{\substack{\rotatebox{-90}{$\in$}\\\Gamma V}}{x})=\lambda\pi_1(x) \tag*{(58 bis)}
\]

Ceci dit, \deleted{\unsure{le fibré|la filt}} tensorisant (47) par $L^{-1}$
on trouve sur $L^{-1}\otimes V$ une structure
d'extension
\[
\unsure{\cdot|0}\longrightarrow\underline O_S\xrightarrow{i'}\underbrace{L^{-1}\otimes V}_{\substack{\wr\,\downarrow\,\alpha\\\check V}}\xrightarrow{\unsure{\pi|\pi'}}\deleted{\unsure{\mathcal O_S}}\check L\longrightarrow0 \tag*{(59)}
\]
% slanted
\margin{Généraliser cet énoncé au cas de la suite exacte (49) et de sa transposée\ldots}
qui, compte tenu de l'iso (56), n'est autre
que la structure d'extension (51) [i.e.
$\alpha$ transforme $i'_1$ en $i'$, $\pi'_1$ en $\pi'$].

Considérons l'hom composé
\diagram{\begin{tikzcd}[ampersand replacement=\&,column sep=small,row sep=small]
\& \mathrm{Sym}^n(\check V)\arrow[r,hook] \& \mathrm{Sym}^{*}(\check V)=\mathrm{Alg\,aff}\,\check{\mathbb V}(V)\arrow[r] \& \mathrm{Sym}^{*}(\check V)/(e_V-1)\arrow[d,equal]\\
\& L^{\otimes(-n)}\underline{\mathrm{Sym}}^n(V)\arrow[u,"\sim"',"\mathrm{Sym}^n(\alpha)=\alpha_n"]\arrow[d,"\sim"',"\mathrm{Sym}^n(\alpha)=\alpha_n"] \& \& \mathrm{Alg\,aff}(F)\\
\text{soit}\quad\text{(60)} \& \underline{\mathrm{Sym}}^n(\check V)\arrow[r,"j_n"]\arrow[d,equal] \& \mathrm{Alg\,aff}(F) \& \\
\& \substack{\text{faisceau des}\\\text{$n$-formes sur }V} \& \&
\end{tikzcd}}
% archive sheet number: 36
